\BourbakiExercises{习题}

\begin{enumerate}
\def\labelenumi{\arabic{enumi})}
\tightlist
\item
  \begin{enumerate}
  \def\labelenumii{\alph{enumii})}
  \tightlist
  \item
    设 D 是一个以 K 为中心的非交换体，其元素在 K 上的次数都 \(\leqslant 2\)。证明 D 是 K 上的一个四元数代数（VIII, 第 270 页，习题 6, b)）。
  \end{enumerate}
\end{enumerate}

\begin{enumerate}
\def\labelenumi{\alph{enumi})}
\setcounter{enumi}{1}
\tightlist
\item
  证明一个是以 K 为中心的非交换体的凯莱 K-代数是 K 上的一个四元数代数。
\end{enumerate}

¶ 2) a) 设 A 是一个在其中心 K 上有限秩的单环，s 是 A 的一个对合反自同构。设 \(s'\) 是 A 的一个在 K 上与 s 一致的对合反自同构。证明存在 A 的一个元素 a，使得我们有 \(s'(x) = as(x)a^{-1}\)（对每个 \(x \in A\)）且 \(s(a) = a\) 或 \(s(a) = -a\)（应用斯科伦--诺特定理与 V, §11, No.~6, 第 85 页，定理 3）。也证明其逆。

\begin{enumerate}
\def\labelenumi{\alph{enumi})}
\setcounter{enumi}{1}
\item
  设 s 在 K 上的限制不是恒等映射；设 \(K_{0}\) 是 K 中被 s 固定的元素所成的子域。证明集 \(A_{0}\)（相应地，\(A_{1}\)）——由 A 中满足 \(s(a) = a\)（相应地，\(s(a) = -a\)）的元素 a 组成——是 A 的一个 \(K_{0}\) -线性子空间，其维数为 {[}A : K{]}，且 \(A_{0}\) 是 A 上的一个 \(K_{0}\) -结构（用 V, §10, No.~4, 第 63 页，命题 7）。
\item
  设 s 在 K 上的限制是恒等映射；令 \([A:K]=m^{2}\)。证明 \(A_{0}\) 是 A 的一个 \(K_{0}\) -线性子空间，其维数为 \(m(m+1)/2\) 或 \(m(m-1)/2\)（化归为 A 是 K 上矩阵代数的情形并用 a)）。
\end{enumerate}

¶ 3) 设 \((F, s)\) 是一个半单（结合）凯莱 K-代数。

\begin{enumerate}
\def\labelenumi{\alph{enumi})}
\item
  证明 K-代数 F 是单的，或同构于一个单代数与它的反代数的乘积（考虑共轭在中心幂等元集上的作用）。
\item
  证明：若 F 不是交换的，则它同构于 K 上的一个四元数代数（用习题 1）。
\item
  证明：若 F 是交换的，则我们处于下列情形之一：
\end{enumerate}

\(\alpha)\) F = K，且 s 是恒等映射。

\(\beta)\) F = K × K，且 \(s(\lambda, \mu) = (\mu, \lambda)\)。

\(\gamma)\) F 是 K 的一个可分二次扩张，且 \(s\) 是 F 的非平凡自同构。

\(\delta)\) K 的特征为 2，F 是 K 的一个高度为 1 的 2-根式扩张，且 \(s\) 是恒等映射。

¶ 4) 设 \((F, s)\) 是一个阿廷（结合）凯莱 K-代数，\(\mathfrak{r}\) 是它的根。a) 证明我们有 \(x^{2} = 0\) 与 \(s(x) = -x\)（对每个 \(x \in \mathfrak{r}\)）；代数 \(\overline{\mathrm{F}} = \mathrm{F} / \mathfrak{r}\) 是凯莱的且同构于习题 3 中描述的代数之一。

\begin{enumerate}
\def\labelenumi{\alph{enumi})}
\setcounter{enumi}{1}
\item
  证明：若 \(\overline{F}\) 是 K 上的一个四元数代数，则 \(\mathfrak{r}\) 约化为 0（注意我们有 \((ab - ba)x = 0\)（对 \(a, b \in F, x \in \mathfrak{r}\)））。
\item
  设 \(\overline{F}\) 是交换的。证明 \(\mathfrak{r}^{3}\) 约化为 0，但 \(\mathfrak{r}^{2}\) 未必如此。
\item
  设 \(\overline{F}\) 是交换的且 \(\mathfrak{r}^{2}\) 约化为 0。证明存在一个 \(\overline{F}\) -模 V，使得 F 同构于 \(\overline{F} \oplus V\)，且它被赋予使 \((\lambda, v)(\mu, w) = (\lambda \mu, \lambda w + s(\mu) v)\) 成立的 K-代数结构。（在习题 3, c) 的情形 \(\alpha\) ) 至 \(\gamma\) ) 中，用 VIII, 第 245 页，推论 1。在情形 \(\delta\) ) 中，选取 \((e_{i})_{i \in I}\)，它是 \(\overline{F}\) 在 K 上的一个 p-基，并对每个 \(i \in I\)，选取一个元素 \(f_{i}\)（属于 F）提升 \(e_{i}\)，然后证明由 \((f_{i})\) 生成的 F 的子代数同构于 \(\overline{F}\)。
\end{enumerate}

¶ 5) 设 K 的特征 ≠ 2。设 F 是一个不一定结合的交错凯莱 K-代数（III, 附录， No.~2）；我们把它的共轭记作 \(s : x \mapsto \overline{x}\)，把它的凯莱迹与凯莱范数分别记作 T 与 N。我们说 F 的一个子代数 E 是非退化的，如果对 E 的每个非零元素 x，存在一个 \(x' \in E\) 使得 \(\mathrm{T}(x\overline{x}') \neq 0\)。

\begin{enumerate}
\def\labelenumi{\alph{enumi})}
\item
  证明等式 \(x(\overline{z}y) + z(\overline{x}y) = \mathrm{T}(\overline{x}z)y = (yx)\overline{z} + (yz)\overline{x}\)（对 F 中的 x, y, z；化归为 z = x 的情形）。
\item
  设 E 是 F 的一个包含 1、异于 F、在 s 下稳定且在 K 上有限维的子代数；设 E 与 F 都是非退化的。证明存在 F 的一个非零元素 j，使得 \(T(xj) = 0\)（对每个 \(x \in E\)）且 \(N(j) \neq 0\)。我们有 \(\overline{j} = -j\)，\(xj = -j\overline{x}\)（对每个 \(x \in E\)），以及 \(j^{2} = \gamma1_{F}\)（其中 \(\gamma = -N(j)\)）。
\item
  设 \(x\) 与 \(y\) 是 E 的元素。证明公式 \(x(yj) = (yx)j = (yj)\overline{x}\) 与 \((xj)(yj) = \gamma \overline{y}x\)（用 \(a\) ) 与 III, 附录， 第 654 页，习题 2）。
\end{enumerate}

\(d\) ) 证明 \(\mathrm{E} + \mathrm{E}j\) 是 F 的一个非退化子代数，它在 \(s\) 下稳定，维数为 \(2\dim (\mathrm{E})\)，且可等同于 \((\mathrm{E},s)\) 的、由 \(\gamma\) 定义的凯莱扩张。由 III, 附录， No.~2, 第 614 页，命题 3 推出 E 是结合的。

\begin{enumerate}
\def\labelenumi{\alph{enumi})}
\setcounter{enumi}{4}
\tightlist
\item
  由上述结果推出，非退化交错凯莱 K-代数是下列这些：
\end{enumerate}

\begin{enumerate}
\def\labelenumi{(\roman{enumi})}
\item
  赋予恒等对合的 K-代数 \(\mathrm{K}\)
\item
  K 上型为 \((\alpha,0)\)（其中 \(\alpha\neq0\)）的二次代数
\item
  K 上型为 \((\alpha,0,\gamma)\)（其中 \(\alpha\gamma\neq0\)）的四元数代数
\item
  \(\mathrm{K}\) 上型为 \((\alpha, 0, \gamma, \delta)\)（其中 \(\alpha \gamma \delta \neq 0\)）的八元数代数。
\end{enumerate}

（构造一个凯莱扩张的序列 \(K \subset E_{1} \subset E_{2} \subset \cdots \subset F\)，并注意一个八元数代数不是结合的。）

\begin{enumerate}
\def\labelenumi{\alph{enumi})}
\setcounter{enumi}{5}
\tightlist
\item
  证明一个其非零元素都可逆的交错凯莱 K-代数是上述类型之一。
\end{enumerate}

\(g\) ) 设 R 是一个极大序域。证明一个其非零元素都可逆的交错凯莱 R-代数同构于 R、C、H 或 O（VIII, 第 368 页，注 1）。

\begin{enumerate}
\def\labelenumi{\arabic{enumi})}
\setcounter{enumi}{5}
\tightlist
\item
  设 H 是 Q 上型为 \((-1,0,-1)\) 的四元数体。
\end{enumerate}

\begin{enumerate}
\def\labelenumi{\alph{enumi})}
\item
  \(\mathrm{K}\)（\(\mathbf{Q}\) 的一个扩张）分裂 \(\mathrm{H}\)，当且仅当 \(-1\) 是 \(\mathrm{K}\) 中两个平方之和（应用 VIII, 第 363 页，命题 3 与 VIII, 第 364 页，命题 4）。
\item
  设 K 是 Q 的一个次数为 \(2^{n}\) 且分裂 H 的循环扩张。证明 K 的任何真子域都不分裂 H（注意 K 的次数为 \(2^{n-1}\) 的唯一子域是 K 与一个包含 Q 的极大序域的交）。
\item
  设 k 是一个整数 \(\geqslant1\)，p 是 \(2^{2^{k}}+1\) 的一个素因子。证明使得 \(2^{n}\) 整除 p-1 的最大整数 n 是 \textgreater k（注意我们有 \(2^{p-1}\equiv1\)（模 p））。设 \(\Omega\) 是 1 的一个本原 p-次根，\(\mathrm{E}=\mathbf{Q}(\Omega)\)。证明 -1 是 E 中两个平方之和（设 F 是扩张 \(\mathrm{E}(\sqrt{-1})\)；证明 -1 属于 \(\mathrm{N}_{\mathrm{F}/\mathrm{E}}(\mathrm{F}^{*})\)，办法是证明这对 \(\omega\) 与 \(1+\omega^{h}\) 成立并利用恒等式 \(\prod_{j=0}^{r-1}(1+\mathrm{T}^{2^{j}})=\sum_{h=0}^{2^{r}-1}\mathrm{T}^{h})\)。
\end{enumerate}

\(d\) ) 证明 E 中所含的 \(\mathbf{Q}\) 的次数为 \(2^{n}\) 的循环扩张分裂 H（应用 VIII, 第 283 页，推论 2）。

¶ 7) 设 K 的特征异于 2。

\begin{enumerate}
\def\labelenumi{\alph{enumi})}
\item
  设 D 是一个包含 K 的体，F 是 K 上型为 \((\alpha, 0, \gamma)\) 的四元数代数。证明 \(D \otimes_{K} F\) 不是一个体，当且仅当存在 D 中两个交换的元素 \(x\) 与 \(y\) 使得 \(x^2 - \alpha y^2 = \gamma\)（把 VIII, 第 351 页，习题 4, c) 先应用于 \(D \otimes_{K} K(i)\)，再应用于 \(D \otimes_{K} F\)）。
\item
  设 \(F'\) 是 K 上的另一个四元数代数。则 \(F \otimes_{K} F'\) 是一个体，当且仅当 \(\mathrm{N}_{\mathrm{F}}(q) = \mathrm{N}_{\mathrm{F}'}(q')\)（其中 \(q \in F\) 与 \(q' \in F'\) 是纯四元数（III, §2, 第 618 页，习题 3））的唯一解是 \(q = q' = 0\)。（化归为 F 是一个体的情形并用 a) 的判别法。注意若 x 与 y 都不属于 K，则 y 属于 \(\mathrm{K}(x)\)，且写 \(x = s + q\)（其中 \(x \in K\)，q 是一个纯四元数）。）
\end{enumerate}

\begin{enumerate}
\def\labelenumi{\arabic{enumi})}
\setcounter{enumi}{7}
\tightlist
\item
  设 \(K_{0}\) 是一个特征异于 2 的交换域，\((\mathrm{X}_{n})_{n\geqslant1}\) 与 \((\mathrm{Y}_{n})_{n\geqslant1}\) 是两个无限变量序列，K 是有理分式域 \(\mathrm{K}_{0}((\mathrm{X}_{n}),(\mathrm{Y}_{n}))\)。对任何 \(n\geqslant1\)，用 \(F_{n}\) 表示 K 上型为 \((\mathrm{X}_{n},0,\mathrm{Y}_{n})\) 的四元数体。
\end{enumerate}

\begin{enumerate}
\def\labelenumi{\alph{enumi})}
\item
  证明张量积 \(F = \bigotimes_{n} F_{n}\)（III, §4, No.~5, 第 470 页）是一个以 K 为中心的体，使得由有限子集生成的每个子域在 K 上都有有限秩（用习题 7, a)）。
\item
  证明：若 u 是 F 的一个内自同构，则存在 N 的一个有限子集 J，使得 u 在 \(F_{n}\) 上的限制（对 \(n \notin J\)）是恒等映射。由此推出，存在 F 的非内自同构的一个不可数集。
\end{enumerate}

